%! TEX root = ../main.tex

\chapter{Tabulated Data}
\label{app:tables}

\begin{table}[h!]
\begin{minipage}[t]{0.32\linewidth}
	\caption{Attenuation function parameters. Details in 
	\cref{sssec:reconstruction-energy-estimator}, 
	\cite{aabMeasurementCosmicrayEnergy2020}}
	\label{tab:cic-parameters}
	\begin{tabularx}{\linewidth}{@{\extracolsep{\fill}} cccc}
  	\toprule
	\hline
	& $a$ & $b$ & $c$ \\
	\hline
	$I_1$ & $-0.9$ & $-1.64$ & $0.952$ \\
	$I_2$ & $-0.04$ & $-0.42$ & $0.06$ \\
	$I_3$ & $1.3$ & $0.09$ & $-0.37$ \\
  	\bottomrule
	\end{tabularx}
\end{minipage}\hfill
\begin{minipage}[t]{0.60\linewidth}
	\caption{\SD energy scale calibration parameters for different zenith 
	bins. $N$ is the number of selected shower events used for 
        cross-calibration. From \cite{aabMeasurementCosmicrayEnergy2020}.}
	\label{tab:energy-cross-calibration}
	\begin{tabularx}{\linewidth}{@{\extracolsep{\fill}} cccc}
  	\toprule
	\hline
        & $N$ & $A\,/\,\SI{1e17}{\eV}$ & $B$ \\
	\hline
        $\SI{0}{\degree} < \theta < \SI{30}{\degree}$ & $435$ & $1.89\pm0.08$ & 
	$1.029\pm0.012$ \\
        $\SI{30}{\degree} < \theta < \SI{45}{\degree}$ & $1641$ & $1.86\pm0.04$ 
	& $1.030\pm0.006$ \\
        $\SI{45}{\degree} < \theta < \SI{60}{\degree}$ & $1226$ & $1.83\pm0.04$ 
	& $1.034\pm0.006$ \\
  	\bottomrule
	\end{tabularx}
\end{minipage}%
\end{table}
